\begin{textsample}{POS Dim 1 – pa – Score 48.00 – pa/pa\_em\_82.txt}  \label{ex:f1_pos_037}
d ) itinerâncias da Área de Ciências Humanas e Sociais Aplicadas : A área de Ciências Humanas e Sociais Aplicadas apresentou na nucleação da formação geral básica \textbf{um} conjunto de \textbf{categorias} necessárias para ampliar o repertório sociocultural na formação da juventude paraense, diretamente integradas à nucleação voltada para o mundo do trabalho, no intuito de consolidar os saberes e ressignificar práticas educativas com vistas a aprofundar processos de \textbf{ensino-aprendizagem} comprometidos com a formação humana integral, a partir da sua \textbf{correlação} com os eixos estruturantes da formação para o mundo do trabalho.
Neste \textbf{tópico}, retomam -se as \textbf{problematizações} \textbf{que} constam na área de CHSA e \textbf{que} se correlacionam com o mundo do trabalho.
Tal nucleação, concebe a \textbf{abordagem} das itinerâncias a partir de quatro eixos estruturantes, a saber : investigação científica, processos criativos, mediação e intervenção sociocultural e empreendedorismo social.
Assim, é necessário estabelecer os possíveis elos entre os princípios \textbf{norteadores} da educação básica, as \textbf{categorias} de área e os eixos estruturantes ( CHSA ), e ao buscar suas \textbf{correlações}, compreender a \textbf{inter-relação} desses elementos importantes para garantir \textbf{uma} formação humana integral.
As \textbf{categorias} específicas da área Tempo e Espaço ; Território e \textbf{Fronteira} ; Natureza e Cultura ; Sociedade, Indivíduo, Identidade e Interculturalidade ; Ética, Política e Trabalho retornam nesta nucleação com o objetivo de aprofundar e ampliar os objetos de conhecimentos, \textbf{que} \textbf{uma} vez consolidados na formação geral básica se reorganizam e a partir dessa \textbf{inter-relação} com os eixos estruturantes estabelecem possibilidades de aprendizagem flexível, integrada e omnilateral a partir das unidades curriculares da nucleação.
Assim a \textbf{categoria} de área Tempo e espaço, elucidada no texto da formação geral básica, permite \textbf{que} o olhar do \textbf{sujeito} se direcione para a historicidade da existência humana a partir de \textbf{um} dado contexto local, ao possibilitar \textbf{uma} prática educativa integrada, contextualizada, interdisciplinar e intercultural em contextos diversos da Amazônia paraense.
Nesta perspectiva, evidencia -se \textbf{que} tal \textbf{categoria} é fundamental no envolvimento dos quatro eixos estruturantes : \textbf{dependendo} da intencionalidade pedagógica, \textbf{um} deles poderá ser mais notório \textbf{que} outro, por exemplo, ao estabelecer os procedimentos básicos da pesquisa científica, torna -se \textbf{basilar} para o proceder científico, pois situa o objeto de conhecimento, condição essencial para a compreensão de temáticas a serem aprofundadas a partir dos recortes temporais e espaciais no processo de compreensão profunda de fenômenos sociais, culturais, ambientais.
Portanto, confere -se materialidade do eixo investigação científica, articulador dos demais.
Por outro \textbf{lado}, a importância da construção de diagnósticos da realidade a partir da metodologia da pesquisa social, possibilita a partir da articulação entre o local e global, afinar -se ao eixo mediação e intervenção sociocultural, ao \textbf{qualificar} a ação política dos \textbf{sujeitos}, no seu processo de interpretar o mundo a partir dos conceitos básicos dos saberes da área, desenvolver ações fundamentadas em epistemologias questionadoras das hegemonias dominantes, sendo capaz de construir alternativas contra-hegemônicas \textbf{baseadas} na sustentabilidade, na criticidade, na democracia, enfim na emancipação do \textbf{sujeito}.
Com efeito, a \textbf{categoria} Território e \textbf{fronteira} aprofunda a \textbf{categoria} precedente, porém, entram em \textbf{cena} as lutas pela terra e por poder, as quais estão \textbf{permeadas} pelo jogo político e pela configuração de instituições e organizações \textbf{que} tratam de controlar os processos históricos, as mobilizações e os fenômenos sociais/territoriais, bem como tentam obter legitimidade em suas atuações.
Assim, os eixos estruturantes dialogam com essas \textbf{problematizações} \textbf{que} poderão contribuir para a formação de \textbf{um} olhar epistemológico dos diferentes territórios, inclusive ao conceber a escola \textbf{enquanto} território educativo vinculado a \textbf{uma} comunidade ( CHARLOT, 2014 ), a qual, por sua vez, também é \textbf{um} espaço de variadas \textbf{territorialidades} \textbf{que} compõe o tecido social dos diferentes sujeitos do Ensino Médio.
Esses, ao ressignificarem as suas leituras de mundo, poderão integrar seus projetos de vida na sua dimensão social através de \textbf{um} diagnóstico dessa realidade, bem como da apropriação de ferramentas necessárias para \textbf{uma} maior mobilização e engajamento social, fruição de processos criativos e outras condições importantes para transformações históricas desse lugar ; articula -se, portanto, os quatros eixos novamente a partir do aprofundamento dessa \textbf{categoria}.
O mesmo \textbf{raciocínio} pode ser aplicado à terceira \textbf{categoria} de área ( CHSA ) : Natureza e cultura, \textbf{que} promove questionamentos de \textbf{suma} importância para a compreensão do mundo, ao desnaturalizar determinados processos históricos e fenômenos sociais, evidenciando \textbf{que} os projetos de existência(s) humana(s) sobre a Terra são construídos a partir do conhecimento \textbf{que} o ser humano desenvolve sobre o seu habitat e o modo como ele tenta se apropriar da natureza ; e ao moldá -la de acordo com suas necessidades e superficialidades, provoca \textbf{inúmeros} conflitos e movimentos culturais em prol ou contra ele mesmo.
Desse modo, é evidente \textbf{que} a ação educativa, ao abordar essa \textbf{categoria}, deve também estar conectada aos quatro eixos estruturantes do mundo do trabalho, seja por meio dos projetos integrados da área ou da seleção de campos \textbf{que} permitam \textbf{um} aprofundamento de estudos a partir das necessidades apresentadas em processo de escuta permanente da comunidade escolar.
No \textbf{que} se refere à investigação científica, a \textbf{categoria} supracitada é igualmente importante como reunião de objetos de conhecimento \textbf{que} podem servir de possibilidades de pesquisa para os jovens estudantes, pois no contexto amazônico muitas questões relacionadas às dinâmicas envolvendo os \textbf{homens} e as mulheres com a natureza são significativas e devem fundamentar a integração de saberes e práticas da escola e sua comunidade, afinando -se, por sua vez, ao eixo dos processos criativos.
Esse eixo oferece \textbf{um} olhar crítico às manifestações culturais, para além do seu aspecto folclórico, ao compreendê -las dentro de \textbf{um} conjunto de representações simbólicas, de natureza material e imaterial, produzidas \textbf{historicamente} pela ação humana, através do trabalho \textbf{enquanto} transformação da natureza.
Esse agir dos diferentes \textbf{sujeitos}, pensados em perspectiva de \textbf{uma} leitura social da cultura, oferece ferramentas para a mobilização social do estudante, quando este se compreende \textbf{enquanto} produto e sujeito da cultura.
Assim, contemplamos, a partir dessa \textbf{categoria}, os eixos mediação e intervenção sociocultural e o empreendedorismo social, pois a luta pela preservação do patrimônio histórico e cultural torna -se condição fundamental ao exercício da cidadania plena, e a escola \textbf{enquanto} espaço de reflexão contribuirá à valorização dos bens culturais de sua comunidade.
Nesse aspecto, a \textbf{categoria} Sociedade, Indivíduo, Identidade e Interculturalidade é também \textbf{basilar} para a área de CHSA, na medida em \textbf{que} promove a reflexão profunda da nossa vida social e trata de questões \textbf{identitárias}, as quais envolvem \textbf{inúmeros} aspectos históricos, geográficos, sociológicos, filosóficos e \textbf{antropológicos}.
Nesse sentido, foram estabelecidos vários campos de saberes e práticas eletivos, os quais são fundamentais para \textbf{que} se amplie a visão crítica a respeito de tão complexa questão.
Outro ponto \textbf{imprescindível} nessa discussão é o conceito de interculturalidade, na medida em \textbf{que}, no estado do Pará, coabitam \textbf{inúmeros} povos, com todo \textbf{um} arcabouço cultural próprio, povos com sua ancestralidade e \textbf{que} precisam ser compreendidos em \textbf{um} contexto muito mais amplo ; para isso foram idealizados campos de saberes e práticas eletivos, bem como o projeto de vida e os projetos integrados de ensino, \textbf{que} possibilitam transcender o olhar e gerar empatia com respeito à diversidade cultural, de modo a desconstruir preconceitos oriundos do senso comum, consolidando o exercício da alteridade, como perspectiva para conceber modelos diferentes de \textbf{humanidades}. \textbf{Outrossim}, no processo de formação cidadã, no \textbf{que} se refere ao empreendedorismo social, verifica -se a importância de entender na prática o exercício dos direitos humanos, principalmente o direito à diferença, e aprofundar o debate ético e a compreensão sobre o ativismo cidadão e político.
Nesse sentido, o eixo investigação científica é importante e, \textbf{certamente}, dialoga com essa \textbf{categoria}, na medida em \textbf{que} o estudante deve aprender a manusear \textbf{métodos} de análise para melhor compreender a sua realidade e o seu entorno.
Por outro \textbf{lado}, a pesquisa também poderá contemplar processos criativos, ao identificar e ampliar habilidades do pensar e fazer criativos dos múltiplos indivíduos \textbf{que} interagem em diferentes contextos da Amazônia paraense.
Já os eixos mediação e intervenção sociocultural e empreendedorismo social se conectam a essa \textbf{categoria} quando \textbf{qualificam} as práticas sociais e culturais, por meio de propostas de intervenção \textbf{que} estimulem os jovens a \textbf{uma} vivência mais efetiva com as problemáticas discutidas na nucleação da formação geral básica, seja em relação às questões territoriais, os desrespeitos à diversidade, bem como ao potencial econômico das comunidades tradicionais, por exemplo.
Portanto, esses eixos estruturantes estabelecem \textbf{uma} vivência dos conceitos e dos debates realizados na nucleação da formação geral básica, para \textbf{que} na nucleação da formação para o mundo do trabalho possam não só ser aprofundados, mas também ampliados e diversificados.
A quinta e última \textbf{categoria} de área Ética, Política e Trabalho é indispensável à área de CHSA, para sustentar aspectos fundamentais desta nucleação.
Trata -se, portanto, de \textbf{um} grande mobilizador \textbf{conceitual} para a formação integral do ser humano em sua omnilateralidade.
Atualmente, os debates em torno das questões ético-políticas estão em voga, e são temas \textbf{que} precisam ser desdobrados com muita atenção pelo docente, pois se trata de \textbf{abordagens} \textbf{que} trazem consigo muitas polêmicas.
Cabe ao docente, portanto, elevar a discussão a aspectos ontológicos, de modo a deixar de \textbf{lado} preconceitos relacionados a tais conceitos, iniciando os \textbf{discentes} na investigação científica, nos processos criativos, na mediação e na intervenção sociocultural, bem como consolidando \textbf{uma} base \textbf{conceitual} robusta para \textbf{que} o jovem se compreenda \textbf{amplamente}, percebendo \textbf{que} pode ser \textbf{um} empreendedor social e/ou \textbf{um} trabalhador assalariado comprometido com a sua realidade social, \textbf{que} sabe se posicionar \textbf{enquanto} \textbf{ator} político \textbf{que} luta por igualdade social, por seus direitos, por justiça social e com respeito aos seus semelhantes.
Desse modo, tem -se o cenário \textbf{ideal} para a construção de \textbf{uma} cidadania autônoma, ativa e crítica.
A \textbf{inter-relação} entre princípios curriculares, competências e habilidades de área, \textbf{categorias} e eixos estruturantes será condição fundamental para \textbf{que} os variados objetos do conhecimento sejam aprofundados pelos estudantes.
Nesse momento tem -se \textbf{uma} tematização específica sobre como o educador deve conceber a formação dos estudantes no \textbf{escopo} da flexibilização curricular, materializada em : projetos integrados da área, campos de saberes e práticas eletivos e projeto de vida.
A idealização e construção dos planos de curso realizados pelos docentes da área de CHSA devem levar em consideração o quadro organizador direcionado para o mundo do trabalho, posto \textbf{que} é através da lógica impressa no referido quadro \textbf{que} o docente pode \textbf{ver} claramente as conexões entre todos os elementos \textbf{que} compõem a referida área ; o mesmo pode ser consultado ao fim desta seção.
Com efeito, os professores precisam estar atentos para não tornar essa parte flexível do currículo em \textbf{mero} apêndice da formação geral básica, pois possuem intencionalidades específicas.
No entanto, não são partes concorrentes ; estão entrelaçadas pelo projeto de vida, seja \textbf{enquanto} unidade curricular, seja como meio fundamental para integralizar a formação dos múltiplos sujeitos da educação básica em sua etapa conclusiva.
A \textbf{tarefa} dos profissionais da educação nesse novo cenário não é fácil, mas ao aguçar o seu planejamento de forma integrada, pode promover processos educativos \textbf{que} lhe permitam ir além do usual e gerar experiências educativas relevantes.
Por isso, as competências e habilidades da área, tanto na formação geral básica, quanto na nucleação para o mundo do trabalho, podem ser percebidas dentro da estrutura pedagógica dos dois projetos integrados propostos pela área de CHSA ; o mesmo ocorre com os componentes eletivos propostos para a referida área do conhecimento.
Assim, pode -se perceber a articulação da formação para o mundo do trabalho de acordo com a mesma lógica \textbf{que} \textbf{embasou} a formação geral básica.
O principal objetivo dessa costura entre \textbf{categorias} de área e eixos estruturantes nos projetos integrados da área é o desenvolvimento de competências específicas e habilidades \textbf{que} possam encaminhar e \textbf{inspirar} os estudantes do ensino médio a lançarem mão de sua criatividade, a ponto de elaborarem e consolidarem seus projetos de vida, os quais serão o reflexo das múltiplas identidades das suas juventudes, da diversidade e das perspectivas empreendedoras \textbf{que} vão inseri -los no mundo do trabalho, da ciência, da tecnologia e da cultura.
Dessa forma, é possível contribuir para \textbf{uma} leitura crítica do contexto amazônico, bem como incorporar novos significados e sentidos aos projetos de vida dos estudantes, por meio da preparação, para o exercício da cidadania plena, do desenvolvimento do protagonismo juvenil, e consequentemente no (re)conhecimento dos jovens do Ensino Médio de sua condição de \textbf{sujeitos} éticos, políticos e \textbf{que} estão vinculados ontologicamente com o trabalho.
Por isso, pode -se afirmar \textbf{que} são \textbf{atores} fundamentais na transformação social e histórica da Amazônia paraense.
Incentivar os \textbf{aprendizes} a pensar criticamente e criativamente, a partir de sua omnilateralidade, a \textbf{totalidade} global e os lugares amazônicos a ponto de os mesmos refletirem e agirem sobre \textbf{uma} determinada realidade política, social, econômica, cultural, ambiental como \textbf{sujeitos} históricos e autônomos, é \textbf{imprescindível} na construção de \textbf{um} mundo melhor e com sentido no \textbf{vivido} e na vida destes jovens do ensino médio do Pará.
Orientá -los como mediadores na elaboração, planejamento e execução de projetos didáticos/pedagógicos, valorizando suas juventudes criativas, é a dimensão elementar para conduzir \textbf{uma} política educacional ao longo do processo de conclusão da educação básica.

% matched lemmas: abordagem, amplamente, antropológico, aprendiz, ator, basear, basilar, categoria, cena, certamente, conceitual, correlação, depender, discente, embasar, enquanto, ensino-aprendizagem, escopo, fronteira, historicamente, homem, humanidade, ideal, identitária, imprescindível, inspirar, inter-relação, inúmero, lado, mero, método, norteador, outrossim, permear, problematização, qualificar, que, raciocínio, sujeito, sumo, tarefa, territorialidade, totalidade, tópico, um, ver, viver
\end{textsample}
